\chapter{Conclusion}
\label{cha:conclusion}

This thesis used computer-assisted graph theory to make progress on two open questions concerning the intersection of longest cycles in graphs.

For RQ\ref{rq1}, this was whether there exists a 2-connected graph of circumference 11 in which no vertex meets every 11-cycle. We performed an exhaustive computational search over several restricted classes of 2-connected graphs, using three independent generation pipelines built around \texttt{geng}, \texttt{GenerateUHG}, and \texttt{plantri}, combined with a program implementing a self-designed algorithm (Algorithm~\ref{alg:find-k-intersection}). The generated graphs allowed for making a definitive claim on all 2-connected graphs of order at most $n$ and maximum degree at most $d$ for the $(n,d)$ pairs $(13,6),(14,5),(15,4),(22,3)$, as well as on the 2-connected planar graphs of order 13. For these graphs, no graph was found which had circumference 11, and no vertex meeting all 11-cycles. RQ\ref{rq1} therefore remains open. However, the search uncovered new graphs whose 11-cycle intersection consists of a single vertex, related to a construction by Carol T.\ Zamfirescu. Building on this, we proved (Theorem~\ref{theorem:single-vertex-intersection}) that such single-vertex-intersection graphs exist for every order $n \ge 13$.

For RQ\ref{rq2}, the question was for which $k$ it is true that if a planar graph with minimum degree 4 has a $k$-cycle in every vertex-deleted subgraph, then the graph has a cycle longer than $k$. We reduced the problem to the same cycle-intersection framework (Theorem~\ref{theorem:intersection-vertex-deletion}), and adapted Algorithm~\ref{alg:find-k-intersection} to compute the $k$-cycle intersections for any circumference $k$ of a graph rather than a fixed $k$, giving us Algorithm~\ref{alg:find-all-k-intersection}. Exhaustive search over 2-connected planar graphs of minimum degree at least 4, extended to order 20 using a supplied graph list of the non-Hamiltonian graphs of this class, found no negative example (Observation~\ref{obs:planar-20-nothing}). In parallel, we generalized the gadget underlying Zamfirescu's construction into a new family $G_{p,\ell,D}$ (Theorem~\ref{theorem:d-extension}) with empty longest-cycle intersection, and a planar, minimum-degree-4 subfamily (Theorem~\ref{theorem:new-vals-rq2}) giving new negative examples of RQ\ref{rq2}.

\section{Future work}
RQ\ref{rq1} itself remains unresolved, and thus the question by van Aardt.\ remains open for $k=11$. Extending the exhaustive search to larger orders and degree bounds would require either substantially more computational resources, or algorithmic and generation-side optimizations. One approach that has not been tried in this thesis is to generated cubic graphs, and inserting a vertex on a subset of edges, creating graphs that are similar to the structure of graphs of the family $G_{p,\ell,D}$, as a heuristic method. Alternatively, an attempt at a proof for the nonexistence of a 2-connected, circumference 11 graph with no vertex meeting every 11-cycle may have to be created.

For RQ\ref{rq2}, it would be valuable to characterize more precisely which values of $k$ and for which orders $n$ the family $G_{p,\ell,D}$ covers, as the question still remains open for some pairs $(k,n)$.
